\reportappendix{А}{Полная validation-матрица основной серии}
\label{app:main}
Таблица~\ref{tab:fullmain} построена из компактного CSV принятого аудита. Все значения — число успехов из 20; named3 сюда не входят. On/off в названии конфигурации относится к обучению, в столбцах — к исполнению сохранённых весов. Классика имеет 20/20 в обоих режимах, а best/last к ней не применяются.
\begin{center}\captionof{table}{Best/last и режим исполнения; validation20}\label{tab:fullmain}
\begin{smalltable}\centering\begin{tabular}{@{}lccccc@{}}\toprule
Конфигурация & Seed & Best on & Best off & Last on & Last off\\\midrule
SAC-on & 0 & 20 & 20 & 20 & 0 \\
SAC-on & 1 & 20 & 0 & 20 & 0 \\
SAC-on & 2 & 20 & 20 & 20 & 20 \\
TQC-on & 0 & 20 & 0 & 20 & 0 \\
TQC-on & 1 & 20 & 0 & 20 & 0 \\
TQC-on & 2 & 20 & 20 & 20 & 0 \\
DQN-on & 0 & 20 & 0 & 11 & 0 \\
DQN-on & 1 & 14 & 0 & 14 & 0 \\
DQN-on & 2 & 20 & 0 & 20 & 0 \\
DDPG-on & 0 & 20 & 0 & 20 & 0 \\
DDPG-on & 1 & 0 & 0 & 0 & 0 \\
DDPG-on & 2 & 0 & 0 & 0 & 0 \\
SAC-off & 0 & 0 & 0 & 0 & 0 \\
SAC-off & 1 & 0 & 0 & 0 & 0 \\
SAC-off & 2 & 0 & 0 & 0 & 0 \\
\bottomrule

\end{tabular}\end{smalltable}\end{center}
Показатель относится к выбранным 20 состояниям и не оценивает вероятность успеха на всём пространстве. Лучшие и последние веса относятся к одному обучению и не являются независимыми повторениями эксперимента. Совпадение показателей не означает совпадения весов. Внешние оценки не изменяют веса.

\reportappendix{Б}{Ключи выбора настройки}
\label{app:selection}
Таблица~\ref{tab:keys} иллюстрирует применение~\eqref{eq:config}. Для формального сравнения использованы полные числа из материала \eref{stage1/selection_audit.json}{«Точные ключи выбора настройки»}; округление ниже служит только чтению. Все компоненты рассчитываются по subset=validation, best в режиме обучения, всем трём seeds. Named исключён.
\begin{table}[htbp]\caption{Ключи Stage 1, округлённые до шести знаков}\label{tab:keys}
\begin{smalltable}\begin{tabularx}{\linewidth}{@{}Xrrrrr@{}}\toprule
Вариант & min $p$ & mean $p$ & min $H_f$ & mean $H_f$ & mean $H_\ell$\\\midrule
DDPG baseline & 0 & 0,333333 & 0 & 2,338500 & 2,338500\\
DDPG lr1e4 & 0 & 0,216667 & 0 & 0,924833 & 0,927167\\
DDPG warmup5000 & 0 & 0,666667 & 0 & 4,705000 & 4,705000\\
SAC-off baseline & 0 & 0 & 0 & 0 & 0,073667\\
SAC-off lr1e4 & 0 & 0 & 0 & 0 & 0,001333\\
SAC-off warmup5000 & 0 & 0,333333 & 0 & 2,011333 & 2,126167\\\bottomrule
\end{tabularx}\end{smalltable}\end{table}

Строгий выигрыш warmup5000 установлен на второй компоненте и не зависит от округления длительностей. Это решение об ограниченной последующей проверке, не признание настройки надёжной. Confirmation сохранило требование: каждый новый best $\geq16/20$, минимум два новых last $\geq16/20$, параметры неизменны, неудачные seeds публикуются. Данные seeds 3–5 отвергли подтверждение обеих конфигураций.


\reportappendix{В}{Электронные материалы и воспроизводимость}
\label{app:repro}
\subsection*{Источники чисел и уровень проверки}
Точные машинные таблицы включены в редактируемый комплект. Это внутренние результаты проекта, а не внешние научные публикации. Таблица~\ref{tab:evidence} показывает, какой источник отвечает за каждый раздел. Исходные данные и их контрольные суммы сохранены. Для семи документов подготовлены явно обозначенные переносимые описания; их происхождение отделено от исторических оригиналов.
\begin{table}[htbp]\caption{Компактные источники электронного приложения}\label{tab:evidence}
\begin{smalltable}\begin{tabularx}{\linewidth}{@{}p{32mm}X@{}}\toprule
Раздел & Материал и назначение\\\midrule
Основная серия & \eref{main/evaluation_summary.csv}{Сводки best/last и режимов исполнения}: best/last, фильтр, раздельные validation/named; \eref{main/validation_history.csv}{История валидации}: периодические оценки; \eref{main/AUDIT_S10_S11.md}{аудит S10/S11}.\\
Stage 1 & \eref{stage1/detailed_comparison.csv}{Сравнение вариантов и seed}, \eref{stage1/selection_audit.json}{Точные ключи выбора настройки}: точные ключи; \eref{stage1/hypothesis_assessment.csv}{Проверка диагностической гипотезы}: проверка раннего механизма.\\
Confirmation & \eref{confirmation/seeds_0_5.csv}{Результаты исходных и новых seed}, \eref{confirmation/confirmation_gate.json}{Критерий и исход подтверждения}, \eref{confirmation/hold_component_summary.csv}{Компоненты удержания}: разложение удержания по компонентам.\\
Robustness & \eref{robustness/summary_controller_group.csv}{Контроллеры и группы стартов}: 98 ячеек; \eref{robustness/summary_overall.csv}{Итоги фиксированной матрицы}: итоги; \eref{robustness/paired_outcomes.csv}{Парные исходы off/on}: пары; \eref{robustness/intervention_summary.csv}{Вмешательства фильтра}: коррекции.\\
Модель и протокол & \eref{protocols/RL_ENV_SPEC.md}{Экспериментальный контракт}, \eref{protocols/CART_SAFETY.md}{Защитный фильтр}, \eref{protocols/STRUCTURED_ROBUSTNESS.md}{Протокол начальных условий}; \eref{configs/states.json}{Точный список стартов}: все 140 стартов.\\\bottomrule
\end{tabularx}\end{smalltable}\end{table}

Полнота и происхождение Stage 1 описаны в \eref{protocols/DEVELOPMENT_STAGE1_AUDIT.md}{аудите этапа подбора}; подтверждение — в \eref{protocols/DEVELOPMENT_CONFIRMATION_AUDIT.md}{аудите новых seed}. Поздняя квитанция TQC seed 1 закрывает технический отказ старой очереди переиспользованием тех же 92 эпизодов. Старый текст основного аудита сохраняет прежний статус; это хронология, не новый научный неуспех.

Уровни свидетельств не взаимозаменяемы. SHA256 выявляет изменение байтов, но не является цифровой подписью автора. Проверка компактных таблиц подтверждает суммы, идентичности и применение правил выбора. Предшествующий полный аудит пересчитывал метрики из сохранённых состояний и переходов, не исполняя политики заново. Настоящий отчёт использует уже принятые результаты и проверяет компактные копии; повторного полного аудита или интегрирования динамики не было.

В материале \eref{robustness/analysis_summary.json}{«Сводка определений и результатов»} поле \texttt{raw\_metric\_recomputation=not\_run} относится к экспорту агрегатов. Оно не означает отсутствия отдельного предшествующего аудита полных журналов. Из компактных CSV нельзя восстановить все траектории. Для повторного пересчёта полных метрик нужны внешние архивы и закреплённый анализатор; большие архивы и буферы опыта не включены в комплект передачи. Отдельно доступны \applink{models/README.md}{шесть закреплённых моделей SAC/TQC для исполнения}; их веса и данные проверки не заменяют полные журналы и не позволяют возобновить прежнее обучение.

\subsection*{Проверяемое происхождение}
Электронное приложение сохраняет происхождение 36 материалов, контрольные суммы и сведения о прежних редакциях. Для семи документов переносимые описания явно отделены от исторических оригиналов. Точные идентификаторы научных данных, принятых видео и версий кода приведены в разделе «Происхождение» репозитория. Научные таблицы, изображения и принятые видео сохранены без изменений.

\subsection*{Воспроизведение отчёта}
Исходники отчёта, библиография и команды сборки включены в электронные материалы. Проверка компактных чисел использует стандартный Python без загрузки моделей и запуска среды. Научные графики воспроизводятся из тех же таблиц; в настоящей редакции они сохранены без изменения. Главная страница репозитория содержит краткие команды, а карта приложений связывает разделы отчёта с таблицами, рисунками, моделями и видео.

\subsection*{Происхождение реализации}
Лабораторная модель и исходная классическая основа заимствованы с указанием источника. SAC, TQC, DDPG и DQN исполнены библиотечными реализациями. Вклад проекта — согласованные интерфейсы, экспериментальный протокол, защитный слой, диагностика, сохранение и анализ результатов. Происхождение заимствований и условия использования перечислены в \applink{THIRD_PARTY_NOTICES.md}{сведениях об источниках}; характер помощи при подготовке материалов раскрыт в \applink{docs/AI_USE.md}{описании применения генеративной модели}.

Фотография на Рисунке~\ref{fig:stand_photo} показывает реальную установку. Схематическое изображение на Рисунке~\ref{fig:stand_illustration} поясняет расположение узлов, а схема на Рисунке~\ref{fig:model} задаёт соглашения расчётной модели. Масштаб иллюстраций не используется для извлечения измерений. Анимации построены из сохранённых журналов; эти изображения и ролики не свидетельствуют об аппаратных испытаниях обученных политик.

\reportappendix{Г}{Видеопримеры сохранённых траекторий}
\label{app:videos}
Шесть роликов представлены на странице «Видео экспериментов» электронного приложения. Они созданы из принятых журналов без нового расчёта динамики и имеют фиксированное правило выбора. Таблица~\ref{tab:videos} отделяет иллюстрации от статистики серий. Особенно важно, что название шестого файла не означает проведения закрытого итогового теста.
\begin{table}[htbp]\caption{Видеоприложения: содержание и границы выводов}\label{tab:videos}
\begin{smalltable}\begin{tabularx}{\linewidth}{@{}p{36mm}X@{}}\toprule
Материал & Сценарий и назначение\\\midrule
\vref{01_nominal_controllers}{1. Исходная задача} & Классика, SAC seed 0, TQC seed 0; on, общий nominal\_00. Один выбранный успешный старт.\\
\vref{02_ddpg_best_vs_last}{2. DDPG: лучшие и последние веса} & Confirmation warmup5000 seed 5, validation\_1000, on. Разные веса: 90 и 100 тыс. переходов; пример ограничения скорости при положении маятника около верха.\\
\vref{03_filter_safety_rescue}{3. Предотвращение остановки} & SAC seed 0, angle\_01: off останавливается на 0,25 м, on завершает горизонт с успехом. Первая подходящая пара в закреплённом порядке.\\
\vref{04_success_vs_safety}{4. Успех и граница} & SAC seed 0, angle\_16: off успешен с превышением 0,24 м, on неуспешен без превышения. Один из семи исходов «успех только off».\\
\vref{05_robustness_groups}{5. Семь групп стартов} & SAC seed 0, on; индекс 00 каждой группы. Все выбранные примеры успешны; это не доля успеха группы.\\
\vref{06_final_controllers}{6. Семь контроллеров} & On, общий joint\_00. Три SAC, три TQC и классика; все выбранные примеры успешны. Этот пример не относится к final5000.\\\bottomrule
\end{tabularx}\end{smalltable}\end{table}

Полный manifest видео связывает файлы с эпизодами и исходниками рендера. Сохраняются фактические времена, включая укороченные переходы; после завершения панель застывает. Экранная интерполяция не используется для расчёта метрик. В частности, краткое вмешательство может оказаться между кадрами с частотой 30 кадров/с. Успех и превышение желаемой границы показаны раздельно.

Ролики доступны через главную страницу электронных материалов. Постеры, сведения об отборе, контрольные суммы и принятая проверка качества сохранены. Повторное построение видео требует внешних исходных архивов и старых журналов; они не нужны для просмотра готовых файлов.

\subsection*{Электронные материалы работы}
Код и конфигурации, таблицы и графики, шесть видео, выбранные модели SAC/TQC, протоколы и инструкции воспроизведения доступны через главную страницу репозитория. Репозиторий закрытый (PRIVATE); необходимы авторизация и предоставленный владельцем доступ.
\begin{center}\small\url{\SubmissionRepositoryURL}\end{center}
